% --- Archon physics formalization source begin ---
% archon:physics
% archon:covers PhyXMiniProblems/problem_phyx_mini_0362.lean
% archon:source-report reports/phyx_mini/problem_phyx_mini_0362.source.json

\chapter{Physics problem phyx\_mini\_0362}
\label{ch:PhyXMiniProblems_problem_phyx_mini_0362}

\paragraph{Problem source.}
\#\# Physical scenario

A gas with an initial tempreture of 900\textbackslash{}circ C\textbackslash{} undergoes the process shown in the figure.

\#\# Question

What is the final temperature?

\#\# Answer choices

- A: A: 3246\textasciicircum{}\textbackslash{}circ C

- B: B: 1856\textasciicircum{}\textbackslash{}circ C

- C: C: 1596\textasciicircum{}\textbackslash{}circ C

- D: D: 2556\textasciicircum{}\textbackslash{}circ C

\#\# Figure caption (auxiliary; use the image as primary evidence)

The image is a graph depicting a pressure-volume (p-V) diagram. The x-axis is labeled "V (cm³)" and ranges from 0 to 300, while the y-axis is labeled "p (atm)" and ranges from 0 to 3. The graph shows a horizontal line at 2 atm, indicating a constant pressure process from volume 100 cm³ to 200 cm³. This line is marked with two points: "1" at 100 cm³ and "2" at 200 cm³, with an arrow pointing from point 1 to point 2, indicating the direction of the process.

\#\# Dataset metadata

Ideal Gases and Kinetic Theory; Physical Model Grounding Reasoning, Multi-Formula Reasoning

\paragraph{Recorded answer/context.}
A: A: 3246\textasciicircum{}\textbackslash{}circ C

\paragraph{Figure/image path.}
/root/proposal\_for\_physic/science-mango/phyx\_mini\_run/phyx\_data/test\_image/362.png

\paragraph{Formalization target.}
create a compiling Lean file with sorry bodies at `PhyXMiniProblems/problem\_phyx\_mini\_0362.lean`. The Lean declarations must preserve the physical quantities, dimensions or dimensional roles, figure labels, governing-law hypotheses, and final relation expressed by this problem.
Use Mathlib/Physlib names found through LeanExplore where available. If a physics API is missing, introduce faithful local abstractions rather than scalar placeholder aliases.

\begin{theorem}[Physics formalization target]
\label{thm:physics:phyx_mini_0362:target}
\lean{PhyXMiniProblems.ProblemPhyXMini0362.finalTemperature_matches_recordedAnswerA}
\uses{def:physics:phyx-mini-0362:phyxminiproblems-problemphyxmini0362-temperatureindegreescelsius, def:physics:phyx-mini-0362:phyxminiproblems-problemphyxmini0362-idealgasexpansionsetup, def:physics:phyx-mini-0362:phyxminiproblems-problemphyxmini0362-matchesproblemstatement, def:physics:phyx-mini-0362:phyxminiproblems-problemphyxmini0362-matchesprimarypressurevolumediagram, def:physics:phyx-mini-0362:phyxminiproblems-problemphyxmini0362-hasphysicalidealgasparameters, def:physics:phyx-mini-0362:phyxminiproblems-problemphyxmini0362-satisfiesfixedsampleidealgaslaws, def:physics:phyx-mini-0362:phyxminiproblems-problemphyxmini0362-recordedanswerchoice}
The assigned autoformalize agent should translate this physics problem into Lean declarations in the covered file, with theorem and lemma proof bodies written as `by sorry`.
\end{theorem}
\begin{proof}
This is an autoformalization task, not a proof task. Produce statements that can later be proved by the physics prover without weakening the physical model.
\end{proof}

% --- Archon declaration topology begin ---
\section{Declaration topology}
\begin{definition}[Volume Quantity]
  \label{def:physics:phyx-mini-0362:phyxminiproblems-problemphyxmini0362-volumequantity}
  \lean{PhyXMiniProblems.ProblemPhyXMini0362.VolumeQuantity}
  Physical gas volume, carrying dimension L³ and a nonnegative value
\end{definition}
\begin{proof}
  This is the declared model definition of Volume Quantity.
\end{proof}
\begin{definition}[Pressure Quantity]
  \label{def:physics:phyx-mini-0362:phyxminiproblems-problemphyxmini0362-pressurequantity}
  \lean{PhyXMiniProblems.ProblemPhyXMini0362.PressureQuantity}
  Physical gas pressure, represented by Physlib's dimensionful type
\end{definition}
\begin{proof}
  This is the declared model definition of Pressure Quantity.
\end{proof}
\begin{definition}[Temperature Quantity]
  \label{def:physics:phyx-mini-0362:phyxminiproblems-problemphyxmini0362-temperaturequantity}
  \lean{PhyXMiniProblems.ProblemPhyXMini0362.TemperatureQuantity}
  Absolute thermodynamic temperature
\end{definition}
\begin{proof}
  This is the declared model definition of Temperature Quantity.
\end{proof}
\begin{definition}[pressure In Pascals]
  \label{def:physics:phyx-mini-0362:phyxminiproblems-problemphyxmini0362-pressureinpascals}
  \lean{PhyXMiniProblems.ProblemPhyXMini0362.pressureInPascals}
  \uses{def:physics:phyx-mini-0362:phyxminiproblems-problemphyxmini0362-pressurequantity}
  Coherent-SI pressure readout in pascals
\end{definition}
\begin{proof}
  This is the declared model definition of pressure In Pascals.
\end{proof}
\begin{definition}[pressure In Atmospheres]
  \label{def:physics:phyx-mini-0362:phyxminiproblems-problemphyxmini0362-pressureinatmospheres}
  \lean{PhyXMiniProblems.ProblemPhyXMini0362.pressureInAtmospheres}
  \uses{def:physics:phyx-mini-0362:phyxminiproblems-problemphyxmini0362-pressurequantity, def:physics:phyx-mini-0362:phyxminiproblems-problemphyxmini0362-pressureinpascals}
  Pressure readout in standard atmospheres
\end{definition}
\begin{proof}
  This is the declared model definition of pressure In Atmospheres.
\end{proof}
\begin{definition}[volume In Cubic Centimeters]
  \label{def:physics:phyx-mini-0362:phyxminiproblems-problemphyxmini0362-volumeincubiccentimeters}
  \lean{PhyXMiniProblems.ProblemPhyXMini0362.volumeInCubicCentimeters}
  \uses{def:physics:phyx-mini-0362:phyxminiproblems-problemphyxmini0362-volumequantity}
  Volume readout in cubic centimetres; 1 m³ = 10⁶ cm³
\end{definition}
\begin{proof}
  This is the declared model definition of volume In Cubic Centimeters.
\end{proof}
\begin{definition}[temperature In Kelvins]
  \label{def:physics:phyx-mini-0362:phyxminiproblems-problemphyxmini0362-temperatureinkelvins}
  \lean{PhyXMiniProblems.ProblemPhyXMini0362.temperatureInKelvins}
  \uses{def:physics:phyx-mini-0362:phyxminiproblems-problemphyxmini0362-temperaturequantity}
  Read a stored absolute temperature in kelvins. Physlib's Temperature stores an absolute nonnegative value, while the ratio of temperature units converts the storage scale to TemperatureUnit.kelvin
\end{definition}
\begin{proof}
  This is the declared model definition of temperature In Kelvins.
\end{proof}
\begin{definition}[textbook Celsius Offset]
  \label{def:physics:phyx-mini-0362:phyxminiproblems-problemphyxmini0362-textbookcelsiusoffset}
  \lean{PhyXMiniProblems.ProblemPhyXMini0362.textbookCelsiusOffset}
  The whole-degree Celsius-to-kelvin offset used by the supplied textbook answer choices. This convention is made explicit because using 273.15 would give 3246.30 °C rather than the displayed whole-number answer
\end{definition}
\begin{proof}
  This is the declared model definition of textbook Celsius Offset.
\end{proof}
\begin{definition}[temperature In Degrees Celsius]
  \label{def:physics:phyx-mini-0362:phyxminiproblems-problemphyxmini0362-temperatureindegreescelsius}
  \lean{PhyXMiniProblems.ProblemPhyXMini0362.temperatureInDegreesCelsius}
  \uses{def:physics:phyx-mini-0362:phyxminiproblems-problemphyxmini0362-temperaturequantity, def:physics:phyx-mini-0362:phyxminiproblems-problemphyxmini0362-temperatureinkelvins, def:physics:phyx-mini-0362:phyxminiproblems-problemphyxmini0362-textbookcelsiusoffset}
  Celsius readout under the whole-degree textbook convention
\end{definition}
\begin{proof}
  This is the declared model definition of temperature In Degrees Celsius.
\end{proof}
\begin{definition}[Diagram Point]
  \label{def:physics:phyx-mini-0362:phyxminiproblems-problemphyxmini0362-diagrampoint}
  \lean{PhyXMiniProblems.ProblemPhyXMini0362.DiagramPoint}
  The two numerical labels printed beside the process endpoints
\end{definition}
\begin{proof}
  This is the declared model definition of Diagram Point.
\end{proof}
\begin{definition}[Axis Quantity]
  \label{def:physics:phyx-mini-0362:phyxminiproblems-problemphyxmini0362-axisquantity}
  \lean{PhyXMiniProblems.ProblemPhyXMini0362.AxisQuantity}
  Physical quantities assigned to the two graph axes
\end{definition}
\begin{proof}
  This is the declared model definition of Axis Quantity.
\end{proof}
\begin{definition}[Volume Display Unit]
  \label{def:physics:phyx-mini-0362:phyxminiproblems-problemphyxmini0362-volumedisplayunit}
  \lean{PhyXMiniProblems.ProblemPhyXMini0362.VolumeDisplayUnit}
  Unit printed on the horizontal axis
\end{definition}
\begin{proof}
  This is the declared model definition of Volume Display Unit.
\end{proof}
\begin{definition}[Pressure Display Unit]
  \label{def:physics:phyx-mini-0362:phyxminiproblems-problemphyxmini0362-pressuredisplayunit}
  \lean{PhyXMiniProblems.ProblemPhyXMini0362.PressureDisplayUnit}
  Unit printed on the vertical axis
\end{definition}
\begin{proof}
  This is the declared model definition of Pressure Display Unit.
\end{proof}
\begin{definition}[Segment Shape]
  \label{def:physics:phyx-mini-0362:phyxminiproblems-problemphyxmini0362-segmentshape}
  \lean{PhyXMiniProblems.ProblemPhyXMini0362.SegmentShape}
  Qualitative shape of the directed segment in the bitmap
\end{definition}
\begin{proof}
  This is the declared model definition of Segment Shape.
\end{proof}
\begin{definition}[Process Kind]
  \label{def:physics:phyx-mini-0362:phyxminiproblems-problemphyxmini0362-processkind}
  \lean{PhyXMiniProblems.ProblemPhyXMini0362.ProcessKind}
  Thermodynamic process classification
\end{definition}
\begin{proof}
  This is the declared model definition of Process Kind.
\end{proof}
\begin{definition}[Equation Of State Model]
  \label{def:physics:phyx-mini-0362:phyxminiproblems-problemphyxmini0362-equationofstatemodel}
  \lean{PhyXMiniProblems.ProblemPhyXMini0362.EquationOfStateModel}
  Equation-of-state model for the gas sample
\end{definition}
\begin{proof}
  This is the declared model definition of Equation Of State Model.
\end{proof}
\begin{definition}[Thermodynamic State]
  \label{def:physics:phyx-mini-0362:phyxminiproblems-problemphyxmini0362-thermodynamicstate}
  \lean{PhyXMiniProblems.ProblemPhyXMini0362.ThermodynamicState}
  \uses{def:physics:phyx-mini-0362:phyxminiproblems-problemphyxmini0362-volumequantity, def:physics:phyx-mini-0362:phyxminiproblems-problemphyxmini0362-pressurequantity, def:physics:phyx-mini-0362:phyxminiproblems-problemphyxmini0362-temperaturequantity}
  Pressure, volume, and absolute temperature at one endpoint
\end{definition}
\begin{proof}
  This is the declared model definition of Thermodynamic State.
\end{proof}
\begin{definition}[Gas Sample]
  \label{def:physics:phyx-mini-0362:phyxminiproblems-problemphyxmini0362-gassample}
  \lean{PhyXMiniProblems.ProblemPhyXMini0362.GasSample}
  \uses{def:physics:phyx-mini-0362:phyxminiproblems-problemphyxmini0362-equationofstatemodel}
  A fixed gas sample. The numerical nR field is a readout in atm·cm³/K; it is not used as a replacement type for the gas itself
\end{definition}
\begin{proof}
  This is the declared model definition of Gas Sample.
\end{proof}
\begin{definition}[Pressure Volume Diagram]
  \label{def:physics:phyx-mini-0362:phyxminiproblems-problemphyxmini0362-pressurevolumediagram}
  \lean{PhyXMiniProblems.ProblemPhyXMini0362.PressureVolumeDiagram}
  \uses{def:physics:phyx-mini-0362:phyxminiproblems-problemphyxmini0362-diagrampoint, def:physics:phyx-mini-0362:phyxminiproblems-problemphyxmini0362-axisquantity, def:physics:phyx-mini-0362:phyxminiproblems-problemphyxmini0362-volumedisplayunit, def:physics:phyx-mini-0362:phyxminiproblems-problemphyxmini0362-pressuredisplayunit, def:physics:phyx-mini-0362:phyxminiproblems-problemphyxmini0362-segmentshape}
  Raw labels, coordinate scales, point coordinates, and arrow from the image
\end{definition}
\begin{proof}
  This is the declared model definition of Pressure Volume Diagram.
\end{proof}
\begin{definition}[Ideal Gas Expansion Setup]
  \label{def:physics:phyx-mini-0362:phyxminiproblems-problemphyxmini0362-idealgasexpansionsetup}
  \lean{PhyXMiniProblems.ProblemPhyXMini0362.IdealGasExpansionSetup}
  \uses{def:physics:phyx-mini-0362:phyxminiproblems-problemphyxmini0362-diagrampoint, def:physics:phyx-mini-0362:phyxminiproblems-problemphyxmini0362-processkind, def:physics:phyx-mini-0362:phyxminiproblems-problemphyxmini0362-thermodynamicstate, def:physics:phyx-mini-0362:phyxminiproblems-problemphyxmini0362-gassample, def:physics:phyx-mini-0362:phyxminiproblems-problemphyxmini0362-pressurevolumediagram}
  The fixed gas, its endpoint states, process type, and supplied diagram
\end{definition}
\begin{proof}
  This is the declared model definition of Ideal Gas Expansion Setup.
\end{proof}
\begin{definition}[Matches Problem Statement]
  \label{def:physics:phyx-mini-0362:phyxminiproblems-problemphyxmini0362-matchesproblemstatement}
  \lean{PhyXMiniProblems.ProblemPhyXMini0362.MatchesProblemStatement}
  \uses{def:physics:phyx-mini-0362:phyxminiproblems-problemphyxmini0362-temperatureindegreescelsius, def:physics:phyx-mini-0362:phyxminiproblems-problemphyxmini0362-idealgasexpansionsetup}
  Information stated in the prose. The initial temperature occurs here, while the requested final temperature deliberately does not
\end{definition}
\begin{proof}
  This is the declared model definition of Matches Problem Statement.
\end{proof}
\begin{definition}[Matches Primary Pressure Volume Diagram]
  \label{def:physics:phyx-mini-0362:phyxminiproblems-problemphyxmini0362-matchesprimarypressurevolumediagram}
  \lean{PhyXMiniProblems.ProblemPhyXMini0362.MatchesPrimaryPressureVolumeDiagram}
  \uses{def:physics:phyx-mini-0362:phyxminiproblems-problemphyxmini0362-pressureinatmospheres, def:physics:phyx-mini-0362:phyxminiproblems-problemphyxmini0362-volumeincubiccentimeters, def:physics:phyx-mini-0362:phyxminiproblems-problemphyxmini0362-diagrampoint, def:physics:phyx-mini-0362:phyxminiproblems-problemphyxmini0362-idealgasexpansionsetup}
  Primary-image evidence. The fields record the printed axes and ranges, both point coordinates, the horizontal segment, and its direction. State 2 is at 300 cm³, as shown by the bitmap
\end{definition}
\begin{proof}
  This is the declared model definition of Matches Primary Pressure Volume Diagram.
\end{proof}
\begin{definition}[Has Physical Ideal Gas Parameters]
  \label{def:physics:phyx-mini-0362:phyxminiproblems-problemphyxmini0362-hasphysicalidealgasparameters}
  \lean{PhyXMiniProblems.ProblemPhyXMini0362.HasPhysicalIdealGasParameters}
  \uses{def:physics:phyx-mini-0362:phyxminiproblems-problemphyxmini0362-pressureinatmospheres, def:physics:phyx-mini-0362:phyxminiproblems-problemphyxmini0362-volumeincubiccentimeters, def:physics:phyx-mini-0362:phyxminiproblems-problemphyxmini0362-temperatureinkelvins, def:physics:phyx-mini-0362:phyxminiproblems-problemphyxmini0362-diagrampoint, def:physics:phyx-mini-0362:phyxminiproblems-problemphyxmini0362-idealgasexpansionsetup}
  Positivity conditions selecting physically meaningful gas states
\end{definition}
\begin{proof}
  This is the declared model definition of Has Physical Ideal Gas Parameters.
\end{proof}
\begin{definition}[Satisfies Fixed Sample Ideal Gas Laws]
  \label{def:physics:phyx-mini-0362:phyxminiproblems-problemphyxmini0362-satisfiesfixedsampleidealgaslaws}
  \lean{PhyXMiniProblems.ProblemPhyXMini0362.SatisfiesFixedSampleIdealGasLaws}
  \uses{def:physics:phyx-mini-0362:phyxminiproblems-problemphyxmini0362-pressureinatmospheres, def:physics:phyx-mini-0362:phyxminiproblems-problemphyxmini0362-volumeincubiccentimeters, def:physics:phyx-mini-0362:phyxminiproblems-problemphyxmini0362-temperatureinkelvins, def:physics:phyx-mini-0362:phyxminiproblems-problemphyxmini0362-diagrampoint, def:physics:phyx-mini-0362:phyxminiproblems-problemphyxmini0362-idealgasexpansionsetup}
  The molar ideal-gas equation pV = (nR)T at every labelled state, together with the physical meaning of an isobaric endpoint process. These are general governing relations and contain no numerical final temperature
\end{definition}
\begin{proof}
  This is the declared model definition of Satisfies Fixed Sample Ideal Gas Laws.
\end{proof}
\begin{lemma}[final Absolute Temperature In Kelvins eq 3519]
  \label{lem:physics:phyx-mini-0362:phyxminiproblems-problemphyxmini0362-finalabsolutetemperatureinkelvins-eq-3519}
  \lean{PhyXMiniProblems.ProblemPhyXMini0362.finalAbsoluteTemperatureInKelvins_eq_3519}
  \uses{def:physics:phyx-mini-0362:phyxminiproblems-problemphyxmini0362-temperatureinkelvins, def:physics:phyx-mini-0362:phyxminiproblems-problemphyxmini0362-idealgasexpansionsetup, def:physics:phyx-mini-0362:phyxminiproblems-problemphyxmini0362-matchesproblemstatement, def:physics:phyx-mini-0362:phyxminiproblems-problemphyxmini0362-matchesprimarypressurevolumediagram, def:physics:phyx-mini-0362:phyxminiproblems-problemphyxmini0362-hasphysicalidealgasparameters, def:physics:phyx-mini-0362:phyxminiproblems-problemphyxmini0362-satisfiesfixedsampleidealgaslaws}
  The constant pressure and threefold volume increase make the absolute final temperature 3519 K, since the initial 900 °C is 1173 K under the stated textbook convention
\end{lemma}
\begin{proof}
  The conclusion follows from the governing relations and figure readouts named in the statement of final Absolute Temperature In Kelvins eq 3519.
\end{proof}
\begin{definition}[Answer Choice]
  \label{def:physics:phyx-mini-0362:phyxminiproblems-problemphyxmini0362-answerchoice}
  \lean{PhyXMiniProblems.ProblemPhyXMini0362.AnswerChoice}
  Labels printed beside the four displayed answer choices
\end{definition}
\begin{proof}
  This is the declared model definition of Answer Choice.
\end{proof}
\begin{definition}[celsius Value]
  \label{def:physics:phyx-mini-0362:phyxminiproblems-problemphyxmini0362-answerchoice-celsiusvalue}
  \lean{PhyXMiniProblems.ProblemPhyXMini0362.AnswerChoice.celsiusValue}
  \uses{def:physics:phyx-mini-0362:phyxminiproblems-problemphyxmini0362-answerchoice}
  Celsius value printed beside each answer label
\end{definition}
\begin{proof}
  This is the declared model definition of celsius Value.
\end{proof}
\begin{definition}[recorded Answer Choice]
  \label{def:physics:phyx-mini-0362:phyxminiproblems-problemphyxmini0362-recordedanswerchoice}
  \lean{PhyXMiniProblems.ProblemPhyXMini0362.recordedAnswerChoice}
  \uses{def:physics:phyx-mini-0362:phyxminiproblems-problemphyxmini0362-answerchoice}
  Answer label recorded in the supplied dataset metadata
\end{definition}
\begin{proof}
  This is the declared model definition of recorded Answer Choice.
\end{proof}
% --- Archon declaration topology end ---
% --- Archon physics formalization source end ---
